\documentclass[11pt,a4paper]{article}
\usepackage[margin=1in]{geometry}
\usepackage{amsmath,amssymb,amsthm}
\usepackage{algorithm}
\usepackage{algpseudocode}
\usepackage{booktabs}
\usepackage{hyperref}

\theoremstyle{definition}
\newtheorem{definition}{Definition}[section]
\newtheorem{theorem}{Theorem}[section]
\newtheorem{lemma}[theorem]{Lemma}
\newtheorem{remark}{Remark}[section]

\title{\textbf{Rigorous Mathematical Specification, Dual-Timescale Dynamics, and Convergence of Lookahead and Exponential Moving Average (EMA)}}
\author{Team Scaibu \\ \small Email: \texttt{jaydeep.9664920749@gmail.com} \\ \small Architecture and Core Engine Team}
\date{October 2026}

\begin{document}
\maketitle

\begin{abstract}
This document presents the complete mathematical formulation and algorithmic specification of Lookahead Optimizer dynamics and Exponential Moving Average (EMA) weight smoothing (Polyak-Ruppert averaging). Lookahead decouples optimization across two timescales: an inner fast loop exploring the parameter space with an arbitrary base optimizer (e.g., Adam/SGD), and an outer slow loop that interpolates toward the fast sequence $\phi_{k+1} = \phi_k + \alpha (\theta_{k, J} - \phi_k)$, dampening variance and escaping oscillation traps. We derive variance reduction guarantees, provide the complete algorithmic pseudo-code, detail computational complexities, step through a numerical calculation, and address stability considerations.
\end{abstract}

\section{Notation and Mathematical Preliminaries}

Let $\theta, \phi \in \mathbb{R}^P$ denote parameter vectors representing fast and slow weights respectively.

\begin{itemize}
    \item $\theta_{k, j} \in \mathbb{R}^P$: Fast parameter iterate at outer step $k$ and inner step $j \in \{1, \dots, J\}$.
    \item $\phi_k \in \mathbb{R}^P$: Slow parameter iterate at outer step $k$.
    \item $\alpha \in (0, 1)$: Synchronization interpolation step size (or EMA decay factor $\beta \in (0, 1)$).
    \item $J \in \mathbb{N}_{\ge 1}$: Synchronization period (inner step count, e.g., $J = 5$ or $10$).
    \item $\mathcal{A}$: Arbitrary base optimization algorithm (e.g., AdamW, SGD).
\end{itemize}

\begin{table}[h]
\centering
\caption{Summary of Mathematical Symbols for Lookahead / EMA}
\begin{tabular}{lll}
\toprule
\textbf{Symbol} & \textbf{Type / Domain} & \textbf{Description} \\
\midrule
$\theta$ & $\mathbb{R}^P$ & Fast weight vector \\
$\phi$ & $\mathbb{R}^P$ & Slow / EMA weight vector \\
$\alpha$ & $(0, 1)$ & Interpolation / smoothing rate \\
$J$ & $\mathbb{N}_{\ge 1}$ & Synchronization step interval ($5$) \\
\bottomrule
\end{tabular}
\end{table}

\section{Problem Definition and Core Concepts}

\subsection{Intuition and Motivation}
Standard stochastic optimizers frequently oscillate in high-curvature directions or become trapped in noisy local fluctuations. Lookahead introduces a dual-timescale mechanism: fast weights $\theta$ aggressively search along the loss landscape, while slow weights $\phi$ track a smoothed trajectory. By taking linear steps along the chord connecting $\phi_k$ to $\theta_{k, J}$, Lookahead effectively projects the parameter iterate towards the center of oscillation, drastically reducing variance without requiring hyperparameter retuning. In parallel, continuous Exponential Moving Average (EMA) maintains a running convex combination $\phi_t = \beta \phi_{t-1} + (1 - \beta) \theta_t$, producing superior evaluation checkpoints for inference.

\subsection{Formal Mathematical Formulation}
\begin{definition}[Lookahead and EMA Equations]
\begin{itemize}
    \item \textbf{Lookahead Update} (at every $J$-th inner step):
    \begin{align}
    \phi_{k+1} &= \phi_k + \alpha (\theta_{k, J} - \phi_k) = (1 - \alpha) \phi_k + \alpha \theta_{k, J} \label{eq:lookahead_step} \\
    \theta_{k+1, 0} &= \phi_{k+1} \label{eq:lookahead_reset}
    \end{align}
    \item \textbf{EMA Tracking Update} (per iteration $t$):
    \begin{equation}
    \phi_t^{\text{EMA}} = \alpha \phi_{t-1}^{\text{EMA}} + (1 - \alpha) \theta_t \label{eq:ema_step}
    \end{equation}
\end{itemize}
\end{definition}

\section{Algorithmic Specification}

The computational implementation from \texttt{NnAlgoLookaheadWeightAveraging} is presented below.

\begin{algorithm}[H]
\caption{Lookahead / EMA Weight Averaging Synchronization Step}
\begin{algorithmic}[1]
\Require Fast weights $\theta \in \mathbb{R}^P$, slow weights $\phi \in \mathbb{R}^P$, mode $\in \{\text{lookahead}, \text{ema}\}$
\Require Interpolation / smoothing rate $\alpha \in (0, 1)$
\Ensure Updated slow weights $\phi_{\text{next}} \in \mathbb{R}^P$, synced fast weights $\theta_{\text{next}}$ (for Lookahead)
\State $P \gets \text{length}(\theta)$
\If{$P = 0$ \textbf{or} $\text{length}(\phi) \ne P$}
    \State \textbf{throw} PreconditionFailureException(``Weight buffer lengths mismatch or zero length'')
\EndIf
\If{$\alpha \le 0$ \textbf{or} $\alpha \ge 1.0$}
    \State \textbf{throw} PreconditionFailureException(``Alpha must be strictly in (0, 1)'')
\EndIf
\State $\phi_{\text{next}} \gets \text{new Array of size } P$
\If{$\text{mode} = \text{ema}$}
    \For{$i \gets 1 \textbf{ to } P$}
        \State $\phi_{\text{next}}[i] \gets \alpha \cdot \phi[i] + (1.0 - \alpha) \cdot \theta[i]$
    \EndFor
    \State \Return $\{\text{updated\_slow\_weights}: \phi_{\text{next}}\}$
\Else
    \For{$i \gets 1 \textbf{ to } P$}
        \State $\phi_{\text{next}}[i] \gets \phi[i] + \alpha \cdot (\theta[i] - \phi[i])$
    \EndFor
    \State \Return $\{\text{updated\_slow\_weights}: \phi_{\text{next}}, \text{synced\_fast\_weights}: \phi_{\text{next}}\}$
\EndIf
\end{algorithmic}
\end{algorithm}

\section{Theoretical Guarantees and Variance Reduction}

\begin{theorem}[Variance Reduction on Quadratic Objectives]
Let $f(\theta) = \frac{1}{2} \theta^T A \theta$ be a strongly convex quadratic objective with Hessian $A \succ 0$. Under stochastic gradient updates with gradient covariance $\Sigma$, the asymptotic stationary covariance of Lookahead slow weights satisfies:
\begin{equation}
\operatorname{Cov}(\phi_k) \prec \operatorname{Cov}(\theta_{k, j})
\end{equation}
scaling strictly by factor $\mathcal{O}\left(\frac{\alpha}{2 - \alpha}\right)$.
\end{theorem}

\begin{proof}
Treating inner steps as noisy Markov transitions $\theta_{k, J} = (I - B) \phi_k + e_k$ where $\operatorname{Cov}(e_k) = \Sigma_e$, the outer recurrence becomes $\phi_{k+1} = (I - \alpha B) \phi_k + \alpha e_k$. Solving the Lyapunov discrete covariance equation:
\begin{equation}
\operatorname{Cov}(\phi) = (I - \alpha B) \operatorname{Cov}(\phi) (I - \alpha B)^T + \alpha^2 \Sigma_e
\end{equation}
yielding $\operatorname{Cov}(\phi) \approx \frac{\alpha}{2 B} \Sigma_e$, which decays linearly with $\alpha \to 0$.
\end{proof}

\subsection{Limitations}
\begin{itemize}
    \item \textbf{Synchronization Lag}: If fast weights diverge into an explosive gradient explosion within the $J$ inner steps, the slow weights may still be dragged into the corrupted subspace upon synchronization.
\end{itemize}

\section{Complexity Analysis}

\subsection{Time Complexity}
\begin{itemize}
    \item \textbf{Vector Interpolation}: Executing $P$ linear interpolations takes exactly $P$ multiply-adds ($\mathcal{O}(P)$).
    \item \textbf{Total Time}: $\Theta(P)$ FLOPs per sync event.
\end{itemize}

\subsection{Space Complexity}
\begin{itemize}
    \item \textbf{Auxiliary Memory}: Stores exactly one shadow copy of the parameters $\phi \in \mathbb{R}^P$.
    \item \textbf{Total Space}: $\Theta(P)$ floating point elements ($4P$ bytes in FP32).
\end{itemize}

\section{Worked Numerical Example}

Consider $P=2$, $\theta = [2.0, -1.0]^T$, $\phi = [1.0, 0.0]^T$.

\subsection{Case 1: Lookahead ($\alpha = 0.5$)}
\begin{align}
\phi_{\text{next}, 1} &= 1.0 + 0.5 \times (2.0 - 1.0) = 1.0 + 0.5 = 1.50 \\
\phi_{\text{next}, 2} &= 0.0 + 0.5 \times (-1.0 - 0.0) = 0.0 - 0.5 = -0.50
\end{align}
Fast weights are synchronized to $[1.50, -0.50]^T$.

\subsection{Case 2: EMA ($\alpha = 0.99$)}
\begin{align}
\phi_{\text{next}, 1}^{\text{EMA}} &= 0.99(1.0) + (1 - 0.99)(2.0) = 0.99 + 0.02 = 1.010 \\
\phi_{\text{next}, 2}^{\text{EMA}} &= 0.99(0.0) + 0.01(-1.0) = 0.0 - 0.01 = -0.010
\end{align}

\section{Numerical Considerations and Edge Cases}

\begin{itemize}
    \item \textbf{In-place Buffer Synchronization}: In Lookahead, fast weights must be re-pointed to or cloned from the slow weights buffer to prevent memory leakage.
    \item \textbf{Batch Normalization Statistics}: When evaluating with EMA / Lookahead slow weights, running batch normalization statistics ($\mu, \sigma^2$) should be recalibrated on a subset of training data before validation.
\end{itemize}

\begin{thebibliography}{9}
\bibitem{zhang2019} M.~R.~Zhang, J.~Lucas, G.~Hinton, and J.~Ba, ``Lookahead Optimizer: $k$ steps forward, 1 step back,'' in \emph{Advances in Neural Information Processing Systems (NeurIPS)}, vol.~32, 2019.
\bibitem{polyak1992} B.~T.~Polyak and A.~B.~Juditsky, ``Acceleration of stochastic approximation by averaging,'' \emph{SIAM Journal on Control and Optimization}, vol.~30, no.~4, pp.~838--855, 1992.
\bibitem{izmailov2018} P.~Izmailov et al., ``Averaging weights leads to wider optima and better generalization,'' in \emph{Conference on Uncertainty in Artificial Intelligence (UAI)}, 2018.
\end{thebibliography}

\end{document}
